\section*{Capítulo \ref{ch:b2:lab-techniques-2} --- Técnicas de laboratorio II: síntesis y análisis en la práctica}

\begin{solution}{exo:b2:lab-techniques-2:1}
Hielo y agua; hielo y sal; hielo seco y acetona (como para los \omterm{def:b2:enolates-aldol:enolate}{enolatos} del \cref{ch:b2:enolates-aldol}).
\end{solution}

\begin{solution}{exo:b2:lab-techniques-2:2}
La capa inferior. Añádanse unas gotas de agua: se unen a la capa acuosa, aquí la superior.
\end{solution}

\begin{solution}{exo:b2:lab-techniques-2:3}
Las primeras porciones se apelmazan al captar agua; cuando el sólido recién añadido queda suelto y se mueve libremente
al agitar, la disolución está seca.
\end{solution}

\begin{solution}{exo:b2:lab-techniques-2:4}
El $R_f$ es demasiado alto para una columna (se busca alrededor de 0.3) y las dos manchas están próximas: úsese un eluyente
menos polar, con más hexano (9\,:\,1, por ejemplo), que rebaja ambos $R_f$ y amplía su separación en
volúmenes de columna.
\end{solution}

\begin{solution}{exo:b2:lab-techniques-2:5}
\begin{itemize}
  \item Hidrogenocarbonato de sodio (pH 8.3 $>$ 4.2 + 2): benzoato en el agua; acidúlese y extráigase el ácido benzoico.
  \item Hidróxido de sodio (pH $>$ 12 $>$ 9.99 + 2): fenóxido en el agua; acidúlese y extráigase el fenol.
  \item Ácido clorhídrico (pH de alrededor de 1 $<$ 4.6 $-$ 2): anilinio en el agua; basifíquese y extráigase la anilina.
  \item El naftaleno permanece en el éter: séquese y evapórese.
\end{itemize}
El lavado con hidrogenocarbonato debe ir primero: el hidróxido extraería juntos el ácido y el fenol.
% ledger: ph:NaHCO3, pka:benzoic, pka:phenol, pka:anilinium
\end{solution}

\begin{solution}{exo:b2:lab-techniques-2:6}
$1/T = 1/T_b - (R/\Delta_{\mathrm{vap}}H)\ln(p/p^\circ)$: bromobenceno, unos \qty{49}{\degreeCelsius};
benzaldehído, unos \qty{68}{\degreeCelsius} (figura de la \cref{prop:b2:lab-techniques-2:reduced-pressure}).
\end{solution}

\begin{solution}{exo:b2:lab-techniques-2:7}
$0.050 \times 200\,000/400 = \qty{25}{K}$.
\end{solution}

\begin{solution}{exo:b2:lab-techniques-2:8}
Cada compuesto tiene su propia respuesta: el porcentaje de área solo es un porcentaje en masa si los \omterm{def:b2:chromatography:quantitative}{factores de respuesta} son
iguales. RMN: $0.12/3 = 0.040$~mol de impureza por mol de producto, es decir, $0.040 \times 100 = \qty{4}{g}$
por \qty{200}{g} de producto; $w = 200/204 \approx 98.0~\%$.
\end{solution}

\begin{solution}{exo:b2:lab-techniques-2:9}
Masa teórica $0.0300 \times 164 = \qty{4.92}{g}$. Sin corregir, $4.10/4.92 \approx 83.3~\%$; corregido,
$0.950 \times 4.10/4.92 \approx 79.2~\%$.
\end{solution}

\begin{solution}{exo:b2:lab-techniques-2:10}
El benceno indica que el reactivo se formó y luego se protonó: agua en el material de vidrio, en el disolvente o en el
aldehído; aire (humedad) que entra por una fuga; un aldehído parcialmente oxidado a ácido benzoico, cuyo
hidrógeno ácido destruye el reactivo. Remedios: material secado en estufa, disolvente anhidro, aldehído recién
destilado, una sobrepresión de nitrógeno comprobada en el borboteador. Que casi no se formara reactivo indicaría un
fallo de la iniciación (una superficie de magnesio cubierta de óxido), que se evita con virutas nuevas o activadas.
\end{solution}

\begin{solution}{exo:b2:lab-techniques-2:11}
\qty{270}{kJ} a \qty{150}{W}: como mínimo $270\,000/150 = \qty{1800}{s}$, media hora, y solo si el
haluro reacciona tan deprisa como se añade. Más deprisa, el calor liberado supera la refrigeración: la temperatura sube,
o el reactivo se acumula y puede reaccionar más tarde de golpe (\cref{prop:b2:lab-techniques-2:accumulation}).
\end{solution}

\begin{solution}{exo:b2:lab-techniques-2:12}
Recristalización: $13.0 \times 0.80 = \qty{10.4}{g}$ al 99~\%. Columna: $13.0 \times 0.95 =
\qty{12.4}{g}$ al 99.5~\%, pero con \qty{1.5}{L} de disolvente que comprar, evaporar y eliminar. La columna da
unos \qty{2}{g} más de un producto algo más puro; la recristalización es más sencilla, más barata y genera menos
residuos, y se prefiere cuando su pureza basta.
\end{solution}

\begin{solution}{pb:b2:lab-techniques-2:1}
\textbf{1.} Etoxietano: extremadamente inflamable, nocivo, forma peróxidos. Bromobenceno: inflamable, tóxico,
tóxico para los organismos acuáticos. Magnesio: sólido inflamable, desprende hidrógeno con el agua.
\textbf{2.} \ce{C6H5MgBr + H2O -> C6H6 + Mg(OH)Br}: cada gota de agua destruye reactivo; el secado en estufa
elimina la película de agua del vidrio.
\textbf{3.} Para mantener fuera la humedad y el oxígeno, que también destruye el reactivo.
\textbf{4.} El flujo de nitrógeno; impide que entre aire por la salida.
\textbf{5.} No reacciona con el reactivo, los pares no enlazantes de su oxígeno estabilizan el magnesio, y su
bajo punto de ebullición (\qty{307.7}{K}) permite que el reflujo limite la temperatura.
\textbf{6.} $15.7/157.01 = \qty{0.1000}{mol}$ de bromobenceno; $2.67/24.305 = \qty{0.1099}{mol}$ de magnesio,
en un exceso del 10~\%.
\textbf{7.} $0.1000 \times 270 = \qty{27.0}{kJ}$.
\textbf{8.} Para que el calor se libere al ritmo al que el reflujo y el baño lo evacuan.
\textbf{9.} $0.0200 \times 270 = \qty{5.40}{kJ}$.
\textbf{10.} $5400/170 \approx \qty{32}{K}$.
\textbf{11.} Unos \qty{52}{\degreeCelsius}, muy por encima del punto de ebullición del etoxietano
(\qty{34.6}{\degreeCelsius}): una ebullición violenta que puede desbordar el refrigerante y expulsar vapor inflamable
fuera del matraz.
\textbf{12.} Añádase alrededor de una décima parte del bromobenceno, espérese a que arranque la reacción (turbidez, reflujo
suave), y añádase luego el resto al ritmo de un reflujo suave, con el baño de hielo preparado.
\textbf{13.} \ce{C6H5MgBr + C6H5CHO} da el alcóxido \ce{(C6H5)2CHOMgBr}; el benzaldehído,
$9.55/106.12 = \qty{0.0900}{mol}$, es el limitante.
\textbf{14.} El alcohol es doblemente bencílico: un ácido fuerte lo convertiría en un catión estabilizado,
que daría éteres o productos de sustitución; el cloruro de amonio es lo bastante ácido para protonar el alcóxido
y disolver las sales de magnesio.
\textbf{15.} La capa superior: el etoxietano es menos denso que el agua.
\textbf{16.} Elimina la mayor parte del agua disuelta en el éter.
\textbf{17.} Sulfato de magnesio hasta que quede suelto, filtración, y luego un evaporador rotatorio a presión
reducida.
\textbf{18.} El bromuro de fenilmagnesio se acopla con el bromobenceno sin reaccionar; una concentración local elevada de
bromobenceno y una temperatura alta lo favorecen, otra razón para una adición lenta.
\textbf{19.} El bifenilo ($R_f$ 0.85), el menos polar.
\textbf{20.} Alrededor de $1/0.85 \approx 1.2$ volúmenes de columna para el bifenilo, y $1/0.25 = 4$ para el alcohol.
\textbf{21.} El alcohol aporta 10 protones \omterm{def:b2:huckel:aromatic}{aromáticos} por cada CH; los $10.90 - 10.00 = 0.90$ sobrantes proceden del
bifenilo, con 10 protones por molécula: 0.090 mol de bifenilo por mol de alcohol.
\textbf{22.} $w = 184.24/(184.24 + 0.090 \times 154.21) \approx 0.930$.
\textbf{23.} $12.50 \times 0.930 \approx \qty{11.62}{g}$.
\textbf{24.} $0.0900 \times 184.24 \approx \qty{16.58}{g}$.
\textbf{25.} $12.50/16.58 \approx 75.4~\%$.
\textbf{26.} $11.62/16.58 \approx \boldsymbol{70~\%}$ (70.1~\%), corregido por la pureza.
\end{solution}
